% !TEX root =./ECE444_Practice_main.tex
%
% L15 practice set (Aperture Distributions and Efficiency) -- the continuous
% aperture results that Module 3 samples into arrays from L16 onward.

\fullwidth{\textbf{Learning Objective 3.1: } Describe an aperture distribution,
compute the beamwidth, sidelobe level and aperture efficiency it produces, and
use those results to size an aperture for a specification.
\\[4pt]\normalfont\small Show your work. A \textbf{Documentation} line at the end
is required (write \textbf{None} if you did not collaborate).}

\begin{questions}

\question \textbf{The uniform aperture.} A rectangular horn is fed so that the
field across its $24$\ \text{cm} broad dimension is essentially constant in
amplitude and phase. It operates at $12\ \ghz$.

\begin{parts}

	\part In one or two sentences, state the relationship between the aperture
	distribution and the far-field pattern, and say what the space frequency
	$u = (L/\lambda)\sinp{\theta}$ does for you that the angle $\theta$ does not.
		\begin{solution}[1.2in]
		The far field is the Fourier transform of the aperture distribution: the
		space factor is the integral of the aperture field weighted by the phase
		$e^{jkx\sinp{\theta}}$ across the opening. Written against $u$, the
		pattern shape depends only on the \emph{shape} of the illumination, so
		one curve serves every aperture with that illumination; the size
		$L/\lambda$ only sets how many degrees each unit of $u$ costs.
		\end{solution}

	\begin{minipage}[t]{0.58\linewidth}
		\part Find the aperture length in wavelengths and the half-power
		beamwidth in the broad-dimension plane.
		\begin{solution}[1.3in]
		\eq{ \lambda &= \frac{3\times 10^{8}}{12\times 10^{9}} = 0.025\ \m
		     = 2.5\ \text{cm} \\
		     \frac{L}{\lambda} &= \frac{24}{2.5} = 9.6 \\
		     \theta_{\text{HP}} &= 50.8\mydeg \frac{\lambda}{L}
		     = \frac{50.8\mydeg}{9.6} = 5.29\mydeg }
		\end{solution}
	\end{minipage}\hfill%
	\begin{minipage}[t]{0.37\linewidth}
		\raggedleft \vspace{12pt}
		$\theta_{\text{HP}}$ = \ansbox[1.3in]{$5.29\mydeg$}
	\end{minipage}

	\begin{minipage}[t]{0.58\linewidth}
		\part Find the angle of the first null measured from broadside.
		\begin{solution}[1.1in]
		The sinc pattern nulls at $u = 1$, so
		\eq{ \sinp{\theta} &= \frac{\lambda}{L} = \frac{1}{9.6} = 0.1042 \\
		     \theta &= 5.98\mydeg }
		\end{solution}
	\end{minipage}\hfill%
	\begin{minipage}[t]{0.37\linewidth}
		\raggedleft \vspace{12pt}
		$\theta_{\text{null}}$ = \ansbox[1.3in]{$\pm 5.98\mydeg$}
	\end{minipage}

	\part State the first sidelobe level of this aperture, and explain in one
	sentence why lengthening the horn would not change it.
		\begin{solution}[1.2in]
		The first sidelobe of a uniformly illuminated aperture is $-13.3\ \db$
		relative to the peak. The normalized pattern $\sinp{\pi u}/(\pi u)$
		contains no $L$ at all, so lengthening the aperture rescales the angle
		axis (narrower beam, higher gain) without changing the height of any
		lobe. Sidelobe level is set by the shape of the illumination, and the
		uniform shape is fixed at $-13.3\ \db$ regardless of size.
		\end{solution}

\end{parts}

\newpage

\question \textbf{Aperture efficiency.} Aperture efficiency converts physical
area into gain through $G = \eta_{\text{ap}}\ 4\pi A/\lambda^{2}$.

\begin{parts}

	\part Write the definition of $\eta_{\text{ap}}$ in terms of the aperture
	field, and state in words what the numerator and the denominator each
	represent.
		\begin{solution}[1.3in]
		\eq{ \eta_{\text{ap}} = \frac{\left| \int E_a\ da \right|^{2}}
		     {A \int \left| E_a \right|^{2} da} }
		The numerator is the coherent sum of the aperture field, which is what
		arrives on boresight when every point adds in phase. The denominator is
		the area times the power the aperture actually had to radiate. The ratio
		is coherent gain over available gain, it never exceeds one, and it
		reaches one only for constant amplitude and constant phase.
		\end{solution}

	\begin{minipage}[t]{0.58\linewidth}
		\part A line aperture carries a triangular illumination,
		$E_a = 1 - 2\left| x \right|/L$. Evaluate the two integrals over
		$\xi = x/L$ and compute $\eta_{\text{ap}}$.
		\begin{solution}[1.5in]
		\eq{ \int_{-1/2}^{1/2}\lp 1 - 2\left| \xi \right| \rp d\xi &= \frac{1}{2} \\
		     \int_{-1/2}^{1/2}\lp 1 - 2\left| \xi \right| \rp^{2} d\xi &= \frac{1}{3} \\
		     \eta_{\text{ap}} &= \frac{(1/2)^{2}}{1 \times (1/3)} = 0.750 }
		This is the canonical triangular value, a gain penalty of
		$10\ \mylog{0.75} = -1.25\ \db$ against uniform.
		\end{solution}
	\end{minipage}\hfill%
	\begin{minipage}[t]{0.37\linewidth}
		\raggedleft \vspace{12pt}
		$\eta_{\text{ap}}$ = \ansbox[1.3in]{$0.750$}
	\end{minipage}

	\begin{minipage}[t]{0.58\linewidth}
		\part A pyramidal horn has an $8\ \text{cm} \times 6\ \text{cm}$ mouth
		and a measured aperture efficiency of $0.51$ at $10\ \ghz$. Find its
		gain in dBi.
		\begin{solution}[1.4in]
		\eq{ \lambda &= 0.03\ \m, \qquad A = 0.08 \times 0.06 = 0.0048\ \m^{2} \\
		     \frac{4\pi A}{\lambda^{2}} &= \frac{4\pi (0.0048)}{(0.03)^{2}} = 67.0 \\
		     G &= 0.51 \times 67.0 = 34.2 \\
		     G_{\text{dBi}} &= 10\ \mylog{34.2} = 15.3\ \dbi }
		\end{solution}
	\end{minipage}\hfill%
	\begin{minipage}[t]{0.37\linewidth}
		\raggedleft \vspace{12pt}
		$G$ = \ansbox[1.3in]{$15.3\ \dbi$}
	\end{minipage}

	\begin{minipage}[t]{0.58\linewidth}
		\part Find the effective aperture $A_e$ of that horn, and say in one
		sentence what the difference between $A_e$ and the mouth area
		physically represents.
		\begin{solution}[1.4in]
		\eq{ A_e &= \eta_{\text{ap}} A = 0.51 \times 0.0048 \\
		     &= 0.00245\ \m^{2} = 24.5\ \text{cm}^{2} }
		The missing $23.5\ \text{cm}^{2}$ is area whose field still radiates but
		does not contribute coherently on boresight, mostly because the horn's
		flare puts a quadratic phase error across the mouth and because the
		field tapers toward the edges. None of it is dissipated as heat, which
		is why $\eta_{\text{ap}}$ is not radiation efficiency.
		\end{solution}
	\end{minipage}\hfill%
	\begin{minipage}[t]{0.37\linewidth}
		\raggedleft \vspace{12pt}
		$A_e$ = \ansbox[1.3in]{$0.00245\ \m^{2}$}
	\end{minipage}

\end{parts}

\newpage

\question \textbf{The taper trade.} A line aperture is $40\lambda$ long. The
canonical taper table is reproduced here.

\fullwidth{\centering\small
\begin{tabular}{lcccc}
\hline
Illumination & First sidelobe & HPBW ($\times\ \lambda/L$) & $\eta_{\text{ap}}$ \\
\hline
Uniform      & $-13.3\ \db$ & $0.886$ & $1.00$  \\
Cosine       & $-23\ \db$   & $1.19$  & $0.81$  \\
Triangular   & $-26.5\ \db$ & $1.27$  & $0.75$  \\
Cosine$^{2}$ & $-31.5\ \db$ & $1.44$  & $0.667$ \\
\hline
\end{tabular}}

\begin{parts}

	\begin{minipage}[t]{0.58\linewidth}
		\part Find the half-power beamwidth with uniform illumination.
		\begin{solution}[1.0in]
		\eq{ \theta_{\text{HP}} = \frac{50.8\mydeg}{40} = 1.27\mydeg }
		\end{solution}
	\end{minipage}\hfill%
	\begin{minipage}[t]{0.37\linewidth}
		\raggedleft \vspace{12pt}
		$\theta_{\text{HP}}$ = \ansbox[1.3in]{$1.27\mydeg$}
	\end{minipage}

	\begin{minipage}[t]{0.58\linewidth}
		\part The system requires a first sidelobe no higher than $-25\ \db$.
		Choose the illumination from the table that meets it with the narrowest
		beam, and find the beamwidth you get.
		\begin{solution}[1.3in]
		Cosine gives only $-23\ \db$, so the narrowest-beam choice that meets
		$-25\ \db$ is triangular at $-26.5\ \db$:
		\eq{ \theta_{\text{HP}} &= 1.27 \frac{\lambda}{L}
		     = \frac{1.27 \times 57.3\mydeg}{40} \\
		     &= 1.82\mydeg }
		\end{solution}
	\end{minipage}\hfill%
	\begin{minipage}[t]{0.37\linewidth}
		\raggedleft \vspace{12pt}
		$\theta_{\text{HP}}$ = \ansbox[1.3in]{$1.82\mydeg$}
	\end{minipage}

	\begin{minipage}[t]{0.58\linewidth}
		\part Find the gain penalty in dB that this taper costs relative to
		uniform illumination on the same aperture.
		\begin{solution}[1.0in]
		\eq{ \Delta G = 10\ \mylog{0.75} = -1.25\ \db }
		\end{solution}
	\end{minipage}\hfill%
	\begin{minipage}[t]{0.37\linewidth}
		\raggedleft \vspace{12pt}
		$\Delta G$ = \ansbox[1.3in]{$-1.25\ \db$}
	\end{minipage}

	\begin{minipage}[t]{0.58\linewidth}
		\part Suppose the $1.27\mydeg$ beamwidth is also a requirement. How long
		must the triangular-illuminated aperture be, and what has the sidelobe
		requirement actually cost you? Answer with a number and one sentence.
		\begin{solution}[1.4in]
		\eq{ L^{\prime} = 40\lambda \times \frac{1.27}{0.886} = 57.3\lambda }
		The aperture must be $43\%$ longer to hold the same beamwidth. Low
		sidelobes are bought with aperture, not with cleverness: you either
		accept a wider beam or you build a bigger antenna, and on a real
		platform that length is usually the binding constraint.
		\end{solution}
	\end{minipage}\hfill%
	\begin{minipage}[t]{0.37\linewidth}
		\raggedleft \vspace{12pt}
		$L^{\prime}$ = \ansbox[1.3in]{$57.3\lambda$}
	\end{minipage}

\end{parts}

\newpage

\question \textbf{Scaling a reflector.} A circular reflector of diameter
$D = 1.2\ \m$ is uniformly illuminated to a good approximation and has an
aperture efficiency of $0.60$ at $10\ \ghz$. Use $\theta_{\text{HP}} =
58.4\mydeg\ \lambda/D$.

\begin{parts}

	\begin{minipage}[t]{0.58\linewidth}
		\part Find the half-power beamwidth at $10\ \ghz$.
		\begin{solution}[1.1in]
		\eq{ \lambda &= 0.03\ \m \\
		     \theta_{\text{HP}} &= 58.4\mydeg \frac{0.03}{1.2} = 1.46\mydeg }
		\end{solution}
	\end{minipage}\hfill%
	\begin{minipage}[t]{0.37\linewidth}
		\raggedleft \vspace{12pt}
		$\theta_{\text{HP}}$ = \ansbox[1.3in]{$1.46\mydeg$}
	\end{minipage}

	\begin{minipage}[t]{0.58\linewidth}
		\part Find the gain in dBi at $10\ \ghz$.
		\begin{solution}[1.4in]
		\eq{ A &= \pi \lp \frac{1.2}{2} \rp^{2} = 1.131\ \m^{2} \\
		     \frac{4\pi A}{\lambda^{2}} &= \frac{4\pi (1.131)}{(0.03)^{2}}
		     = 1.579\times 10^{4} \\
		     G &= 0.60 \lp 1.579\times 10^{4} \rp = 9475 \\
		     G_{\text{dBi}} &= 10\ \mylog{9475} = 39.8\ \dbi }
		\end{solution}
	\end{minipage}\hfill%
	\begin{minipage}[t]{0.37\linewidth}
		\raggedleft \vspace{12pt}
		$G$ = \ansbox[1.3in]{$39.8\ \dbi$}
	\end{minipage}

	\begin{minipage}[t]{0.58\linewidth}
		\part The same dish is used at $30\ \ghz$ with the same aperture
		efficiency. Find the new beamwidth and the new gain.
		\begin{solution}[1.4in]
		Tripling the frequency triples $D/\lambda$:
		\eq{ \theta_{\text{HP}} &= \frac{1.46\mydeg}{3} = 0.487\mydeg \\
		     \Delta G &= 20\ \mylog{3} = 9.5\ \db \\
		     G_{\text{dBi}} &= 39.8 + 9.5 = 49.3\ \dbi }
		\end{solution}
	\end{minipage}\hfill%
	\begin{minipage}[t]{0.37\linewidth}
		\raggedleft \vspace{12pt}
		$\theta_{\text{HP}}$ = \ansbox[1.2in]{$0.487\mydeg$} \\[8pt]
		$G$ = \ansbox[1.2in]{$49.3\ \dbi$}
	\end{minipage}

	\part The dish did not change, yet the gain rose by $9.5\ \db$. Explain
	where that gain came from, and name one practical reason a real reflector
	will fall short of the predicted $49.3\ \dbi$ at $30\ \ghz$.
		\begin{solution}[1.4in]
		Gain is $\eta_{\text{ap}}4\pi A/\lambda^{2}$, so it depends on area
		measured in square wavelengths. At $30\ \ghz$ the same $1.131\ \m^{2}$
		spans nine times as many square wavelengths, which is a factor of nine
		or $9.5\ \db$. In practice the aperture efficiency will not hold: the
		surface tolerance and the feed's own illumination taper both scale with
		wavelength, so a surface accurate enough at $10\ \ghz$ introduces
		significant phase error at $30\ \ghz$ and $\eta_{\text{ap}}$ drops.
		\end{solution}

\end{parts}

\newpage

\question \textbf{Reading a measured pattern.} A student measures the pattern of
an unknown aperture antenna at $10\ \ghz$ on the course range. The main lobe is
symmetric about boresight, the first nulls fall at $\pm 3.4\mydeg$, and the first
sidelobes read $19\ \db$ below the peak.

\begin{parts}

	\begin{minipage}[t]{0.58\linewidth}
		\part Assuming the illumination is uniform, find the aperture dimension
		in the measured plane.
		\begin{solution}[1.2in]
		\eq{ \sinp{\theta_{\text{null}}} &= \frac{\lambda}{L} \\
		     L &= \frac{0.03}{\sinp{3.4\mydeg}} = \frac{0.03}{0.0593} \\
		     &= 0.506\ \m = 16.9\lambda }
		\end{solution}
	\end{minipage}\hfill%
	\begin{minipage}[t]{0.37\linewidth}
		\raggedleft \vspace{12pt}
		$L$ = \ansbox[1.3in]{$0.506\ \m$}
	\end{minipage}

	\begin{minipage}[t]{0.58\linewidth}
		\part Predict the half-power beamwidth that this aperture should show.
		\begin{solution}[1.1in]
		\eq{ \theta_{\text{HP}} &= 0.886 \frac{\lambda}{L}
		     = 0.886 \frac{0.03}{0.506} \\
		     &= 0.0525\ \text{rad} = 3.01\mydeg }
		\end{solution}
	\end{minipage}\hfill%
	\begin{minipage}[t]{0.37\linewidth}
		\raggedleft \vspace{12pt}
		$\theta_{\text{HP}}$ = \ansbox[1.3in]{$3.01\mydeg$}
	\end{minipage}

	\part The measured sidelobe is $19\ \db$ down, not $13.3\ \db$ down. Give
	two distinct explanations, one about the antenna and one about the
	measurement.
		\begin{solution}[1.4in]
		The antenna explanation is that the illumination is not uniform. A
		sidelobe between $-13.3\ \db$ and $-23\ \db$ is what a mild amplitude
		taper produces, so the aperture is probably tapered somewhere between
		uniform and cosine, and the beamwidth will be correspondingly wider than
		the $3.01\mydeg$ predicted in part (b). The measurement explanation is
		that the sidelobe was not resolved: too coarse an angular step, a range
		shorter than $2D^{2}/\lambda$, or reflections from the room all fill in
		nulls and can raise or lower a narrow sidelobe by several dB.
		\end{solution}

	\part The same antenna is later found to have $\eta_{\text{ap}} = 0.44$,
	although its amplitude taper alone would predict about $0.81$. Name two
	mechanisms that account for the difference, and state whether each one
	belongs to aperture efficiency or to radiation efficiency.
		\begin{solution}[1.5in]
		Phase error across the aperture and spillover past the aperture edge
		both belong to aperture efficiency: the power leaves the antenna, but
		it does not add coherently on boresight. Feed and strut blockage and
		cross-polarization are two more of the same kind. Conductor and
		dielectric loss in the feed structure belongs to radiation efficiency
		instead, because that power becomes heat and never radiates at all.
		Measured gain carries both, which is why $G = \eta_{\text{rad}} D$ and
		$D = \eta_{\text{ap}}\ 4\pi A/\lambda^{2}$ are separate statements.
		\end{solution}

\end{parts}

\vspace{1.5em}
\noindent\textbf{Documentation:} \rule[-2pt]{0.6\linewidth}{0.4pt}

\end{questions}
